\section{Introduction}
\label{sec:intro}

An Indonesian undergraduate degree ends with the \emph{skripsi}, a thesis whose chapters follow the conventions of the study programme: a methods chapter (\emph{BAB~III}) whose sub-sections are almost fixed in some programmes, a results chapter (\emph{BAB~IV}) whose shape follows the research design, statistics computed in SPSS, Excel or Python, a literature review that must cite real sources, and formal Indonesian that follows the national spelling rules (EYD) and dictionary (KBBI). Students, supervisors and the tools built for them increasingly hand parts of this work to large language models (LLMs). Which model should they use?

The question has become harder, not easier. One vendor SDK now exposes dozens of models from several vendors, at prices per call that differ by two orders of magnitude and often with a reasoning-effort setting; open-weight models run on a laboratory GPU for the price of electricity. Public leaderboards rank models on English knowledge and reasoning~\cite{hendrycks2021mmlu,liang2023helm}, and Indonesian benchmarks measure language understanding, sentiment and school-exam knowledge~\cite{wilie2020indonlu,koto2020indolem,winata2023nusax,koto2023indommlu,wang2024seaeval}. None of them asks whether a model outlines a methods chapter the way an agribusiness or an Islamic family-law programme expects, cites only the sources it was given, lists references that exist, writes an abstract in formal Indonesian without inventing numbers, or computes a correct Mann--Whitney test from a CSV file. Nor do they report what an answer costs or which model produced it---not a trivial question, since a vendor SDK can silently serve a different model than the one requested~\cite{arya2026paperagent}.

We present \emph{AcademicBench}, eight tasks drawn from the skripsi workflow, and use it to compare \NApi{} API models of one SDK, which we grouped into two tiers by name, with \NLocal{} local models under identical Indonesian prompts. Our contributions are:
\begin{itemize}
  \item \textbf{Tasks and items} (Section~\ref{sec:bench}): \NItems{} items: 225 newly written public items and 30 private excerpts of real theses (published as hashes only); deterministic scoring for six tasks and for the deterministic half of the two prose tasks, whose other half comes from LLM judges of other vendors than the candidate's, with leniency calibration and a reliability analysis.
  \item \textbf{Accounting} (Section~\ref{sec:setup}): every call records the model that the SDK reports as answering, the reasoning effort, tokens and the vendor's AI units (converted to US\$); local runs record GPU time (converted at an on-demand price for the same GPU) and measured board energy.
  \item \textbf{Findings} (Section~\ref{sec:results}): the three best API models are not significantly different, while the cost per item of the leader is \CostRatio{} times that of a model \CheapGap{} below it; methodology questions and formal Indonesian are near ceiling for API models, whereas real references, programme-specific outlines and LaTeX repair separate them; most references that Crossref does not know are found in no other database either; on average, local models trail API models on every task, most on references and code and least on formal Indonesian.
  \item \textbf{An open kit}: runner with a cost gate and a budget cap, scorers, sandbox, reference audit, public items, per-item scores and a generator of blind sheets for human raters, in the paper-agent repository (\url{https://github.com/devnolife/paper-agent}).
\end{itemize}
